% Part: first-order-logic
% Chapter: proof-systems
% Section: seqeunt-calculus

\documentclass[../../../include/open-logic-section]{subfiles}

\begin{document}

\iftag{FOL}
      {\olfileid{fol}{prf}{seq}}
      {\olfileid{pl}{prf}{seq}}

\olsection{相继式演算}

虽然许多推导系统以语句的排列为操作对象，但相继式演算以\emph{相继式}为操作对象。
相继式是形如
\[
!A_1, \dots, !A_m \Sequent !B_1, \dots, !B_m,
\]
的表达式，即由相继式符号~$\Sequent$ 分隔的一对语句序列。
两个序列都可以是空的。
相继式演算中的推导是由相继式构成的树，其中最顶端的相继式具有特殊形式（称为“初始相继式”或“公理”），其余每个相继式都由其正上方的相继式通过某条推理规则得出。
推理规则要么处理相继式中的语句（在左边或右边添加、删除或重排它们），要么在规则的结论中引入一个复杂公式。
例如，对于任意的 $\Gamma$、$\Delta$、$!A$ 和~$!B$，$\LeftR{\land}$ 规则允许从 $!A, \Gamma \Sequent \Delta$ 推出 $!A \land !B, \Gamma \Sequent \Delta$，而 $\RightR{\lif}$ 规则允许从 $!A, \Gamma \Sequent \Delta, !B$ 推出 $\Gamma \Sequent \Delta, !A \lif !B$。
（特别地，$\Gamma$ 和~$\Delta$ 可以是空的。）

基于相继式演算的 $\Proves$ 关系定义如下：$\Gamma \Proves !A$ 当且仅当存在某个序列 $\Gamma_0$，使得 $\Gamma_0$ 中的每个 $!A$ 都在~$\Gamma$ 中，并且存在一个以相继式~$\Gamma_0 \Sequent !A$ 为根的推导。
如果相继式~$\Sequent !A$ 有一个推导，则 $!A$ 是相继式演算中的定理。
例如，下面的推导表明 $\Proves (!A \land !B) \lif !A$：

\begin{prooftree}
  \Axiom$!A \fCenter !A$
  \RightLabel{\LeftR{\land}}
  \UnaryInf$!A \land !B \fCenter !A$
  \RightLabel{\RightR{\lif}}
  \UnaryInf$\fCenter (!A \land !B) \lif !A$
\end{prooftree}

在相继式演算中，如果存在 $\Gamma_0 \Sequent$ 的推导（其中每个 $!A \in \Gamma_0$ 都在~$\Gamma$ 中，且相继式的右边为空），则集合 $\Gamma$ 是\emph{不一致的}。
利用规则 \RightR{\Weakening}，任何语句都可以从一个不一致的集合中推导出来。

相继式演算由 Gerhard Gentzen（格哈德·根岑）于20世纪30年代发明。
由于其系统化且对称的设计，它是一种非常有用的形式化方法，可用于发展推导理论。
在相继式演算中寻找推导相对容易，但这些推导往往难以阅读，且它们与证明的联系有时不易看出。
然而，事实证明它是一种非常优雅的处理推导系统的方法，并且许多逻辑都具有相应的相继式演算系统。

\end{document}